\documentclass[11pt]{article}
\usepackage{mathblatt}


\begin{document}
\pruefheftstil{Prozent \textperiodcentered{} Portion 1}{Prüfungsheft P10}
\noindent{\Large\bfseries Prozent \textperiodcentered{} Portion 1}\hfill{\small\color{mbgrau}Prüfungsheft P10}\par\medskip
\pfrueckblick
\pfpaar{\pfkopfzeile{1.}{Grundlage}
$6$ von $15$ Plätzen sind besetzt. Schreibe den Anteil der besetzten Plätze als Bruch, kürze ihn und erweitere auf Hundertstel.
\par\noindent\hfill \leerfeld
\fusshilfe{1: $\tfrac{40}{100}$}}{\pfkopfzeile{2.}{Grundlage}
Schreibe $\tfrac{3}{25}$ als Dezimalzahl.
\par\noindent\hfill \leerfeld
\fusshilfe{2: $0{,}12$}}
\pfabschnitt{Prozent}
\pfstufe[sprozentundanteilumwa]{Prozent und Anteil umwandeln}
\kommtdran{in 2 von 13 Jahren (2020, 2022) \textperiodcentered{} zuletzt 2022 \textperiodcentered{} 1 BE}
\begin{pfaufgabe}{3.}{1 BE \textperiodcentered{} P10 2022 · 1f}
Ein Fünftel aller Jugendlichen erhält kein Taschengeld. Kreuzen Sie den passenden Prozentsatz an:
\par\noindent\hspace*{1.6em}$\square$ 5 \% \quad $\square$ 20 \% \quad $\square$ 25 \% \quad $\square$ 50 \%\par
\fusshilfe{3: 20 \%, Tipp: Prozentsatz}
\end{pfaufgabe}
\pfweiterekopf
\pfpaar{\pfkopfzeile{4.}{eigene Aufgabe}
Schreibe den Bruch $\tfrac{11}{20}$ in Prozent.
\par\noindent\hfill \leerfeld[\%]
\fusshilfe{4: $55\,\%$}}{\pfkopfzeile{5.}{eigene Aufgabe}
Jeder hundertste Besucher bekommt einen Preis. Wie viel Prozent der Besucher bekommen einen Preis?
\par\noindent\hfill \leerfeld[\%]
\fusshilfe{5: $1\,\%$}}
{\small\begin{pfaufgabe}{6.}{1 BE \textperiodcentered{} P10 2020 · 1a}
Ein Förster sagt: „Bei mir wachsen nur 4 \% der gepflanzten Bäume nicht an.“ Kreuzen Sie an, welche Aussage dasselbe bedeutet:
\pfkreuzzeile{Von 10 Bäumen wachsen 4 nicht an.}\pfkreuzzeile{Jeder vierte Baum wächst nicht an.}\pfkreuzzeile{Von 100 Bäumen wachsen 4 nicht an}
\fusshilfe{6: 4 von 100 Bäumen wachsen nicht an, Tipp: Hundertstel}
\end{pfaufgabe}}
\pfpaar{\pfkopfzeile{7.}{eigene Aufgabe nach P10 2020}
$5\,\%$ der Pakete kommen zu spät. Kreuze die Aussage an, die dazu passt.\\ \kreuz{$5$ von $10$ Paketen kommen zu spät.}\\ \kreuz{Jedes 5. Paket kommt zu spät.}\\ \kreuz{$5$ von $100$ Paketen kommen zu spät.}
\fusshilfe{7: $5$ von $100$ Paketen kommen zu spät}}{}
\end{document}